\begin{problem}[第31届全国中学生物理竞赛复赛][电容器与旋转金属板]{22}
如图所示，一电容器由固定在共同导电底座上的 $N+1$ 片对顶双扇形薄金属板和固定在可旋转的导电对称轴上的 N 片对顶双扇形薄金属板组成，所有顶点共轴，轴线与所有板面垂直，两组板面各自在垂直于轴线的平面上的投影重合，板面扇形半径均为 R，圆心角均为 $\theta_0$（$\frac{\pi}{2}\leq\theta_0<\pi$）；固定金属板和可旋转的金属板相间排列，两相邻金属板之间距离均为 s。此电容器的电容 C 值与可旋转金属板的转角 $\theta$ 有关。已知静电力常量为 k。
    \begin{subproblem}
        开始时两组金属板在垂直于轴线的平面上的投影重合，忽略边缘效应，求可旋转金属板的转角为 $\theta$（$-\theta_0 \leq \theta \leq \theta_0$）时电容器的电容 $C(\theta)$；
    \end{subproblem}
    \begin{subproblem}
        当电容器电容接近最大时，与电动势为 E 的电源接通充电（充电过程中保持可旋转金属板的转角不变），稳定后断开电源，求此时电容器极板所带电荷量和驱动可旋转金属板的力矩；
    \end{subproblem}
    \begin{subproblem}
        假设 $\theta_0 = \frac{\pi}{2}$，考虑边缘效应后，第 (1) 问中的 $C(\theta)$ 可视为在其最大值和最小值之间光滑变化的函数
        \begin{equation*}
            C(\theta) = \frac{1}{2}(C_{\max} + C_{\min}) + \frac{1}{2}(C_{\max} - C_{\min})\cos 2\theta
        \end{equation*}
        式中，$C_{\max}$ 可由第(1)问的结果估算，而 $C_{\min}$ 是因边缘效应计入的，它与 $C_{\max}$ 的比值 $\lambda$ 是已知的。若转轴以角速度 $\omega_m$ 匀速转动，且 $\theta = \omega_m t$，在极板间加一交流电压 $V = V_0 \cos \omega t$。试计算电容器在交流电压作用下能量在一个变化周期内的平均值，并给出该平均值取最大值时所对应的 $\omega_m$。
    \end{subproblem}
\end{problem}